% In Default or Not? SSRN working paper source.
% Build: ./build.sh   (figures/tables from figs.py; Tectonic in conda env "tectonic")
% Every number traces to data/interim/late_flags/outcome/FINAL.json (src/late_flags/final.py); see ../NUMBERS_AUDIT.md.
% Legal and form-scope statements trace to src/late_flags/CONTRACTS_AND_SCOPE.md.
\documentclass[11pt]{article}
\usepackage[letterpaper,margin=1.05in]{geometry}
\usepackage[T1]{fontenc}
\usepackage{newpxtext,newpxmath}
\usepackage{microtype}
\usepackage{setspace}\setstretch{1.12}
\usepackage{graphicx}
\usepackage{booktabs,tabularx,threeparttable}
\usepackage{float}
\usepackage[section]{placeins}
\usepackage{enumitem}
\usepackage[font=small,labelfont=bf,labelsep=period]{caption}
\usepackage{xcolor}
\usepackage[round,authoryear]{natbib}
\usepackage[colorlinks=true,linkcolor=black,citecolor=blue!50!black,urlcolor=blue!50!black]{hyperref}
\usepackage{xurl}
\newcommand{\doi}[1]{\href{https://doi.org/#1}{doi:#1}}
\usepackage{titlesec}
\titleformat{\section}{\large\bfseries}{\thesection.}{0.6em}{}
\titleformat{\subsection}{\normalsize\bfseries}{\thesubsection}{0.6em}{}
\setlist{itemsep=2pt,topsep=4pt}
\graphicspath{{figs/}}
\renewcommand{\topfraction}{0.85}\renewcommand{\bottomfraction}{0.6}
\renewcommand{\textfraction}{0.12}\renewcommand{\floatpagefraction}{0.75}
\makeatletter\newcommand{\tabinput}[1]{\@@input #1 }\makeatother

\title{\textbf{In Default or Not?}\\[2pt]
\Large How Funds Report Bonds of Bankrupt Issuers in SEC Form N-PORT}
\author{Rishabh Gyanendra Upadhyay\thanks{Inception42. E-mail:
\href{mailto:uhrishabh@gmail.com}{uhrishabh@gmail.com}.}\\[2pt] \normalsize Inception42}
\date{November 2026}

\begin{document}
\maketitle
\thispagestyle{empty}

\begin{abstract}
\noindent Form N-PORT asks each fund that files it whether each debt security it holds is ``currently in
default'', and the SEC chose to leave the meaning of default to each fund's own methodology. We measure
how far those methodologies diverge for debt of issuers that have filed for bankruptcy. The sample is
57 Chapter~11 and Chapter~7 cases of SEC-reporting issuers from 2019 to 2026, each verified from the
issuer's Form 8-K, with 73\% of holdings in 2020 cases. First, among bankruptcy-linked bond-dates
(the same bond on the same report date) held by funds of at least two different advisers, the advisers'
fund families give opposite answers in 141 of 305 (46.2\%; 95\% interval 35.5\% to 56.0\%, clustered
over the 50 cases involved). Second, at each fund's first report after the petition became public and before any identified plan
effective date, 18.8\% of 3,783 bond holdings are reported as not in default (95\% interval 13.7\% to
23.0\%, clustered by case). Prepetition loans, whose coverage by the item the SEC has not specified,
show 67.2\%; pooled, 24.9\%. In the governing documents of the five largest cases, the petition is an
event of default that accelerates the debt automatically, although that term is not universal.
Holders that answer differently report nearly identical prices, and a typical position is 0.03\% of
fund net assets, so the finding concerns the reliability of a status field used for screening and
aggregation, not misvaluation or investor loss. These are not error rates: the form permits fund
methodologies to differ. A pre-registered run gave 30.3\%; a later source-based audit removed
securities of unrelated companies and non-filing affiliates, and the corrected estimates are post hoc.

\medskip\noindent\textbf{Keywords:} Form N-PORT, default status, bankruptcy, regulatory reporting,
data quality, mutual funds, ETFs.\\
\textbf{JEL classification:} G23, G28, G33.
\end{abstract}

\newpage
\setcounter{page}{1}

%=====================================================================
\section{Introduction}

Item C.9 of Form N-PORT asks a fund, for each debt security it holds, whether the security is
``currently in default'' \citep{sec_nport_form}. The answers are public for the third month of each
fiscal quarter. They feed screens that count defaulted exposure across funds, data products sold to
risk teams, and research panels built from N-PORT \citep{cavani2025}. When the SEC adopted the form, a
commenter warned that the item requires ``a subjective legal analysis on an instrument-by-instrument
basis, on which conclusions among funds may vary''. The Commission declined to define default and
chose ``to allow funds to continue to use their own methodology'' \citep{sec_2016_release}. The form
also lets funds rely on ``their own internal methodologies and the conventions of their service
providers'', provided they are consistently applied \citep{sec_nport_form}.

This note measures how much those methodologies diverge in the case where divergence is easiest to
see: debt of an issuer that has filed a bankruptcy petition. We make two comparisons. The first is
internal and needs no view on what default means: for the same security on the same report date, do
different fund families give the same answer? The second is external: how often is debt of a verified
petitioner reported as not in default? A petition is a well-dated, public event, and in the governing
documents of the five largest cases in our sample it is an event of default that accelerates the debt
automatically (Section~\ref{sec:item}). It is not a legal finding about any fund's report.

The SEC staff has raised the question in comment letters. In 2020 Goldman Sachs Trust told the staff
that ``the respective Form N-PORT filings did not denote all of such securities as being in default''
\citep{corr_gs2020}; in 2021 the staff asked RiverNorth Specialty Finance to report past-due loans as in
default or in arrears \citep{corr_rivernorth2022}. Those letters concern single registrants. We measure
prevalence, persistence and whether reported prices differ.

The findings are descriptive. Fund families disagree on 46\% of the bond-dates that more than one
family holds (95\% interval 35.5\% to 56.0\%; 39.8\% if each case is weighted equally). Around a fifth of bond holdings (18.8\%) and two thirds of loan holdings are reported as not in
default at the first post-petition report. Holders that disagree on the label report nearly the same
prices. The study was pre-registered and run once on 60 cases, which gave 30.3\%; the corrected
sample removes securities that the frozen rules had wrongly linked to two cases, and we report it as a
post-hoc audit (Section~\ref{sec:record}).

%=====================================================================
\section{The form item, the contracts and related work}\label{sec:item}

\paragraph{The item.} Item C.9 is headed ``For debt securities, also provide''; C.9.c asks ``Currently
in default?'' and C.9.d whether ``any interest payments are in arrears or have any coupon payments been
legally deferred by the issuer'' \citep{sec_nport_form}. The form lists ``debt'' and ``loan'' as separate
asset types in Item C.4 and does not say whether loans fall under C.9. We found no form instruction,
staff answer or schema rule on the point: the XML schema allows the C.9 fields on any holding, and
filers populate them for loans. We therefore report loans separately and treat them as exploratory.

\paragraph{The contracts.} We read the governing indenture of a note and, where funds held one, the
credit agreement of a term loan in each of the five largest cases (Table~\ref{tab:contracts}), and in
the Endo case also the notes funds held most. In all of them a voluntary petition by the obligor, or
for Endo by the parent, is an event of default and acceleration is automatic (``without any declaration
or other act'', ``shall automatically become due and payable''). Each instrument matches a CUSIP held in
the sample, and each obligor was itself a debtor in the case, except Endo Luxembourg Finance Company I,
which we could not find named on EDGAR; its instruments accelerate on the parent's filing, which
occurred. Appendix Table~\ref{tab:sources} gives the filings, provisions and debtor sources. In all
five cases the issuer also stated in an 8-K that the filing was an event of default that accelerated
its debt. The term is
not universal: in Valley National Bancorp's 2032 subordinated notes a voluntary filing is an event of
default but acceleration requires a declaration by the trustee or 25\% of holders \citep{valley2022}. We
did not read the documents of every security in the sample. The automatic stay of 11 U.S.C.~\S362(a)
restrains enforcement and collection against the debtor; it does not cure an event of default
\citep{usc362, uscourts_ch11}.

\paragraph{Pending amendments.} The Commission's February 2026 proposal to amend Form N-PORT would
remove the delta of convertible bonds from Item C.9, noting that ``funds may use different
methodologies for calculating delta for convertible bonds, which adds to variability in the reported
information and reduces its utility'' \citep{sec_2026_proposal}. It leaves the default and arrears items
unchanged.

\begin{table}[!htb]
\centering\small
\caption{Bankruptcy terms in governing documents of the five largest cases.}\label{tab:contracts}
\begin{tabularx}{\textwidth}{@{}l X l@{}}
\toprule
Case & Document & Petition accelerates \\
\midrule
Frontier Communications & 8.000\% first-lien notes due 2027, indenture of 15 March 2019; term B-1 loans, 2017 credit agreement & automatically \\
Intelsat Jackson & 8.500\% notes due 2024, indenture of 19 September 2018; term loan credit agreement (2013 conformed text) & automatically \\
Chesapeake Energy & 11.5\% second-lien notes due 2025 and FLLO term loan agreement, both 19 December 2019 & automatically \\
Endo International & 6.125\% notes due 2029 and 2021 term loans, 25 March 2021; Par Pharmaceutical 7.500\% notes due 2027, 2019 (filing by the parent) & automatically \\
Whiting Petroleum & 6.625\% notes due 2026, fifth supplemental indenture of 27 December 2017 & automatically \\
\bottomrule
\end{tabularx}
\end{table}

\paragraph{SEC correspondence.} A full-text search of staff correspondence returns several dozen letters
that mention N-PORT together with defaulted securities. Table~\ref{tab:corr} lists five that we read and
that discuss the default or arrears fields directly. Two acknowledge omitted flags, one ends with the
registrant defending its answer, and two concern the opposite direction.

\begin{table}[!htb]
\centering\small
\caption{SEC staff correspondence on N-PORT default and arrears fields. Dates are the dates of the
registrants' letters; EDGAR may have released them later.}\label{tab:corr}
\begin{tabularx}{\textwidth}{@{}l l X@{}}
\toprule
Date & Registrant & What it says \\
\midrule
2020-08-28 & Goldman Sachs Trust & Defaults footnoted in financial statements; N-PORT ``did not denote all of such securities as being in default''; to be corrected. \\
2021-02-26 & OFS Credit Company & Staff asks why no defaults are reported; registrant: none in default or arrears. \\
2021-10-13 & RiverNorth Specialty Finance & Past-due loans not reported in default or arrears; to be reported in future filings. \\
2020-09-24 & Vertical Capital Income Fund & Loans flagged in default; staff asks for the non-income-producing symbol. \\
2025-09-19 & Yieldstreet Alternative Income Fund & Three securities in default on N-PORT, one marked non-income producing. \\
\bottomrule
\end{tabularx}
\end{table}

\paragraph{Related work.} A security-level default flag is already supplied as reference data: ICE's
N-PORT reporting service lists a debt default flag and an interest-in-arrears and payment-in-kind flag
\citep{ice_nport}, and Bloomberg offers a comparable field. These are inputs to reporting, and their
existence means that answers within a fund family may come from a shared feed. Comparing the same
security across funds on the same date is established for prices: \citet{cici2011} document dispersion
in month-end valuations of identical corporate bonds across mutual funds, \citet{morningstar2021}
repeat the comparison with N-PORT, and Empirasign publishes N-PORT valuation outliers
\citep{empirasign}. The closest population is in \citet{manoel2026}, who matches the bonds of public
Chapter~11 debtors (2008 to 2023) to the mutual funds holding them at filing, from N-PORT or N-CSR
disclosures and CRSP holdings, and asks whether ties between those funds and the debtor's financial
adviser affect TRACE bond returns around the filing; it does not use the N-PORT default item. In a bounded search we did not find a study that measures how fund
answers to the default item differ for the same security, or how they compare with bankruptcy
petitions; we do not claim that none exists.

%=====================================================================
\section{Data and design}\label{sec:data}

\paragraph{Holdings.} We use every SEC Form N-PORT quarterly data set from 2019Q4 to 2026Q3
\citep{sec_nport_data}. A data set holds the filings disseminated in a calendar quarter, with mixed report
dates and with amendments that can supersede filings in earlier data sets. For each fund (series) and
report date we keep the latest version across all data sets. A holding is a fund, security and report
date; a fund with several lots of one security counts as reporting default if any lot does.

\paragraph{Events.} We search EDGAR full text for 8-Ks filed from July 2019 to October 2026 that tag Item
1.03 (bankruptcy or receivership) and mention Chapter~11 or Chapter~7. A case is the first such 8-K of a
registrant. We keep a case only if its Item 1.03 text states that petitions were filed and gives a
petition date between 30 days before and one day after the 8-K filing date, which removes confirmation
and exit filings of earlier cases and spurious tags. The 8-K filing date is when the petition became
public on EDGAR.

\paragraph{Securities.} We link debtor securities by issuer name, legal entity identifier and the
six-character CUSIP issuer prefix, restricted to corporate debt and loans held by some fund on or before
the petition. We exclude debtor-in-possession loans, securitisations, pass-through trusts and
securities that matured before the petition. A post-hoc audit (Section~\ref{sec:record}) checked, security
by security, that the issuer filed in the case.

\paragraph{Measurement.} For each case, security and fund we take the fund's first report date after the
8-K filing date and within 120 days of the petition, before the case's next Item 1.03 8-K and before the
plan effective date where one is identified. Effective dates were found for 39 of the 57 cases; for the
other 18 cases (1,438 holdings) the window rests on the 120-day limit and the next Item 1.03 8-K, and we
report the known-end-date subsample separately. The median gap from the 8-K filing to the report is 44
days. Intervals are 95\% bootstrap intervals that resample cases (2,000 draws).

\paragraph{Scope.} The sample is a set of public-issuer bankruptcies whose debt registered funds held.
It is concentrated in the 2020 cycle (3,174 of 4,331 holdings) and in a few large cases, and it excludes
private-company bankruptcies, which account for much of the leveraged loan market. It is not a
representative sample of corporate defaults.

\begin{table}[!htb]
\centering\small
\caption{Sample construction.}\label{tab:flow}
\begin{tabularx}{\textwidth}{@{}X r r@{}}
\toprule
 & Cases & Holdings \\
\midrule
\tabinput{tables/flow.tex}
\bottomrule
\end{tabularx}
\end{table}

%=====================================================================
\section{Results}

\subsection{Same security, same date}

Table~\ref{tab:cross} compares fund families holding the same security on the same report date. A
fund family is the adviser named in the fund's most recent Form N-CEN (Part C), or the N-PORT registrant
when none is named. For each security-date in the analysis sample we take every fund that reports the
security on that date, not only funds for which it is the first post-petition report. A family's answer
is ``in default'' if all its funds report default, ``not in default'' if none do, and mixed otherwise.
The classes in Table~\ref{tab:cross} are mutually exclusive and add to the total; the memo row
overlaps them.

Of the 305 bond-dates held by at least two families, families give opposite answers in 141 (46.2\%,
95\% interval 35.5\% to 56.0\%, clustered by case). In 157 every family reports default, in six none
does, and in one no family contradicts another but one family's own funds disagree. The 305 bond-dates
come from 50 cases and the 141 from 32. Weighting each case equally gives 39.8\% (28.8\% to 50.3\%);
dropping any one case gives between 43.3\% and 49.0\%, and dropping the five cases with the most
bond-dates gives 46.8\%. The share is 50.0\% in bond-dates of 2020 cases and 39.0\% (41 of 105) in the others. For
loans, families disagree in 32 of 52 loan-dates. Across bonds and loans, 956 of the 1,079 holdings
reported as not in default sit in a security-date where another family reports default.

Figure~\ref{fig:cross}b shows one case, chosen as the bond-date with the most families on the smaller
side. On 30 June 2020, eleven weeks after Frontier Communications filed, 46 fund families reported its
8.000\% first-lien notes due 2027 (CUSIP 35906ABC1): 20 families (43 funds) reported them in default
and 26 families (45 funds) not in default. The median prices were 101.7 and 101.6. The three largest
holdings on each side are in N-PORT filings 0001145549-20-050750, 0001752724-20-163403 and
0001752724-20-174015 (in default) and 0001752724-20-174964, 0001752724-20-178862 and
0000810893-20-005528 (not in default). The indenture of these notes makes a voluntary petition an
event of default with automatic acceleration (Table~\ref{tab:contracts}).

\begin{table}[!htb]
\centering\small
\caption{Security-dates (same CUSIP, same report date) in the analysis sample, by the answers of
different fund families. Post hoc. Intervals: 95\%, bootstrap over cases.}\label{tab:cross}
\begin{tabularx}{\textwidth}{@{}X r r r@{}}
\toprule
 & Bonds & Loans & All \\
\midrule
\tabinput{tables/cross_family.tex}
\bottomrule
\end{tabularx}
\end{table}

\begin{figure}[!htb]
\centering
\includegraphics{fig1_cross.pdf}
\caption{(a) Share of fund families reporting the bond in default, for each bond-date held by at least
two families (a family whose funds disagree counts as one half). (b) The bond-date with the most
families on the smaller side.}
\label{fig:cross}
\end{figure}

This comparison does not depend on how default is read: different families answer the same question
about the same instrument on the same date differently. It does not tell us which answer is right, and
families are not fully independent sources, since several may take the field from the same data
vendor.

\subsection{Comparison with the petition}

Of 3,783 bond holdings, 18.8\% are reported as not in default at the first report after the petition
became public (Table~\ref{tab:est}). For prepetition loans the share is 67.2\%
(548 holdings in 21 cases); because the SEC has not said whether loans fall under the item, we do not
read the loan figure as comparable with the bond figure. Pooled, 1,079 of 4,331 holdings (24.9\%) are
reported as not in default; weighting each case equally gives 27.2\%, and by reported value 34.7\%
(\$2.82 billion of \$8.12 billion). Restricting to the 39 cases with an identified effective date gives
25.8\%. The holdings involve 847 funds in 285 fund families.

\begin{table}[!htb]
\centering\small
\caption{Share of holdings reported as not in default at the first report after the petition.
Intervals: 95\%, bootstrap over cases.}\label{tab:est}
\begin{tabularx}{\textwidth}{@{}X r c r r@{}}
\toprule
 & Share & Interval & Holdings & Cases \\
\midrule
\tabinput{tables/estimates.tex}
\bottomrule
\end{tabularx}
\end{table}

\subsection{Distribution across cases, fund families and years}

The share varies widely by case (Figure~\ref{fig:dist}a). Among the 47 cases with at least five
holdings, the quartiles are 9.7\%, 22.2\% and 40.4\%. Across all 57 cases, every holding is flagged in
14 and none is in two. The five largest cases account for 61\% of unflagged holdings; dropping any one
case leaves between 22.6\% and 26.0\%, and dropping the five largest leaves 21.3\%.

Of the 95 fund families with at least ten holdings, 26 report every holding as in default and none
reports none; the median family leaves 12.5\% unflagged (Figure~\ref{fig:dist}b). The five largest
families hold 45\% of the unflagged holdings, and dropping them leaves 22.0\%. These family patterns are
descriptive: they partly reflect which issuers each family holds. Within the security-dates where
families disagree, of the 117 families that appear in at least five, 83 give each answer at least once,
30 always report default and 4 never do. Most families are therefore not applying a single fixed
convention. Table~\ref{tab:years} gives the estimate by petition year.

\begin{figure}[!htb]
\centering
\includegraphics{fig2_distribution.pdf}
\caption{(a) Pooled share reported as not in default in each bankruptcy case. (b) The same share for
fund families with at least ten holdings in the sample.}
\label{fig:dist}
\end{figure}

\begin{table}[!htb]
\centering\small
\caption{By petition year (pooled bonds and loans).}\label{tab:years}
\begin{tabular}{@{}l r r r@{}}
\toprule
Year & Cases & Holdings & Not in default \\
\midrule
\tabinput{tables/years.tex}
\bottomrule
\end{tabular}
\end{table}

\subsection{What happened next}

We follow each unflagged holding through the fund's later reports until one year after the petition,
the case's next Item 1.03 8-K or an identified effective date, whichever comes first
(Figure~\ref{fig:pers}). Of the 1,079 holdings, 428 have no later report in that window, so their
persistence is not observed, and 156 leave the portfolio (141 sold or exchanged into securities we
cannot link, 13 exchanged into linked securities, 2 in funds that stopped filing). That leaves 495
holdings observed again while still held. At the next report, 229 (46\%) are reported as in default and
266 (54\%) are not; 163 remain unflagged across at least two later reports. By the end of the window,
nine more are flagged at a later report (238 in all), so 257 are still unflagged when last observed.
When a flag appears, the median time from the 8-K filing is 109 days.

\begin{figure}[!htb]
\centering
\includegraphics{fig3_persistence.pdf}
\caption{Status of the unflagged holdings at the end of the observation window (the earliest of one
year after the petition, the next Item 1.03 8-K and an identified effective date). ``Flagged at a later
report'' counts 229 holdings flagged at the next report and nine flagged afterwards.}
\label{fig:pers}
\end{figure}

\subsection{Prices}

In the 168 security-dates where some holders report default and others do not, the median prices of
the two groups differ by a median of 0.0 points (interquartile range $-0.19$ to $+0.21$); in 87.5\% of
these cells the two medians are within one point (Figure~\ref{fig:prices}). The price comparison groups
holders rather than families, so its 168 security-dates are the 173 of class A in
Table~\ref{tab:cross} less six bond-dates in which no holder reporting ``not in default'' gives a usable
price (in four they report zero principal and zero value, in two a quantity in units rather than
principal amount), plus the one class-D bond-date, where funds of a single family disagree. The median price in these
cells is about 35 per 100 of face value. Only one unflagged holding carries forward its own
pre-petition price exactly. Holders that answer the default question differently report nearly the
same values. This does not validate either group's marks, and it does not show which label is right.

\begin{figure}[!htb]
\centering
\includegraphics{fig4_prices.pdf}
\caption{Median price per 100 of face value in the 168 security-dates where some holders report the
security in default and others do not.}
\label{fig:prices}
\end{figure}

\subsection{Arrears and amendments}

Eleven unflagged holdings report interest in arrears; 1,127 flagged holdings do. Arrears or legally
deferred interest need not be a current contractual default, so we report this without interpretation.
Of 89 holdings whose report was amended, none changed its default answer.

%=====================================================================
\section{Interpretation, relevance and limitations}\label{sec:interp}

\paragraph{What the numbers are not.} They are not error rates. The SEC left default to fund
methodology, so a fund can report ``not in default'' under a methodology that, for example, counts only
missed payments, or treats debt as current while adequate-protection interest is paid, or follows a
contract that does not accelerate automatically. We cannot observe which methodology a fund used. Different methodologies,
shared data vendors, operational differences and plain mistakes could all contribute, and the data do
not identify the reason for any particular disagreement. What we measure is how often reported status departs from a verified petition, and how often funds answering
the same question about the same security disagree.

\paragraph{Why it matters.} In the security-dates where holders disagree, the median reported price is
about 35 per 100 of face value and the two groups' prices agree, so these securities were already
heavily discounted when the labels diverged; the data do not show when the discount arose or how much
of it the petition caused. A typical unflagged position is small
relative to its fund (median 0.03\% of net assets; 90th percentile 0.35\%). The \$2.8 billion figure is
reported exposure, not a loss or an overstatement of net asset value. The relevance is for users who
treat the field as a status signal: a screen that counts defaulted holdings across funds or fund
families, a monitor of distressed exposure built from N-PORT, or a research panel that uses the flag
as a variable; a recent Federal Reserve note, for example, reports the share of defaulted holdings of
interval and tender-offer funds from N-PORT \citep{george2026}. For such users a single fund's answer is not a reliable indicator of an issuer event,
and agreement across holders or a check against issuer filings is needed. We do not claim investor
harm.

\paragraph{Loans.} The loan share is three to four times the bond share. With no SEC statement on
whether loans belong in the item, the difference may reflect reporting scope as much as practice, and
we leave it as an exploratory observation.

\paragraph{Sample.} The study plan set a Stage~0 gate of at least 300 externally dated default events.
Independent construction from 8-K evidence produced 60 confirmation cases, so the gate was not met.
Proceeding with the 60 cases was a deviation from the plan, decided by the author and recorded in the
protocol before it was frozen. Clustering by case accounts for dependence among holdings of the same
issuer; the small number of cases limits precision, and the year-by-year estimates outside 2020 rest on
few cases.

%=====================================================================
\section{Research record}\label{sec:record}

\begin{enumerate}
\item \textbf{Development.} The pipeline was built on cases observed in one quarterly data set (2025Q4):
four cases and 265 holdings, of which 36\% were unflagged. These cases were seen before the protocol was
frozen and are excluded from all estimates.
\item \textbf{Frozen confirmation.} The protocol (primary question, sample, estimator, verdict rule) and
the analysis code were hashed before the run. The single run on 60 cases (Figure~\ref{fig:stages}) gave 30.3\% (20.5\% to
39.8\%). A separate re-implementation from the raw N-PORT files, given only the protocol and the case and
security lists, reproduced every count exactly and confirmed ten randomly chosen petition dates.
\item \textbf{Linkage errors found after the run.} The per-case output showed 225 holdings of Gates
Global, an unrelated company, attached to the Frontier Communications case (some funds had labelled
Gates bonds with Frontier's name, and the frozen rules then accepted Gates' CUSIP prefix), and notes of
Ferrellgas, L.P., the operating partnership, attached to the petition of Ferrellgas Partners, L.P., which
did not include it. Both errors raise the share.
\item \textbf{Audit (post hoc).} Two amendments, hashed before any corrected number was computed, defined
a security-level debtor audit. One model labeller classified all 89 issuer groups; a second, independent
reviewer classified every security in the 37 groups holding 92\% of holdings and 97\% of value, and
searched for each case's plan effective date. They disagreed on two securities, resolved from the cited
documents. Securities of other companies (225 holdings), of non-filing affiliates (103) and with
unresolved status (1) were removed, as were holdings reported after an identified effective date.
\item \textbf{Loan review (post hoc).} All 33 loans were checked for debtor-in-possession and roll-up
status from their titles and public documents before their answers were read. One, McDermott's
superpriority term loan, was rolled into the DIP facility and is excluded (32 holdings).
\item \textbf{Review additions (post hoc).} After an external review we added the same-security
comparison across fund families (Table~\ref{tab:cross}), read the governing documents of the five
largest cases (Table~\ref{tab:contracts}), and report bonds as the main estimate and loans as
exploratory. After a second review we made the classes of Table~\ref{tab:cross} mutually exclusive
(the first version listed overlapping rows), added the case-clustered interval and case count for the
same-security comparison, chose the worked example in Figure~\ref{fig:cross}b by a stated rule, and
added the source table for the contracts (Appendix Table~\ref{tab:sources}). None of these changes the
sample or the estimator.
\item \textbf{Reconciliation.} Applying the audit to the re-implementation's own holdings reproduces the
final counts exactly (4,331 holdings, 1,079 unflagged, 57 cases). A separate rebuild of the
same-security comparison, from the re-implementation's holdings and families taken directly from the raw
N-CEN files without the analysis code or its output, reproduces every count in Table~\ref{tab:cross}.
\end{enumerate}

\begin{figure}[!htb]
\centering
\includegraphics{fig5_stages.pdf}
\caption{Pooled share of holdings reported as not in default, by stage of the research record. The
frozen run is the pre-registered confirmation; the later stages are post-hoc corrections of the linked
securities.}
\label{fig:stages}
\end{figure}
The frozen estimate remains part of the record. Because the linkage errors were found after the run,
the audited estimates govern what we conclude, and they are post hoc rather than independently
confirmed.

%=====================================================================
\section{Conclusion}

Fund families holding the same bond on the same date give opposite answers to the N-PORT default
question in nearly half of the bond-dates that more than one family holds, and about a fifth of bond
holdings of bankrupt issuers are reported as not in default at the first report after the petition.
Their reported prices agree. The disagreement occurs in a field for which the SEC permits different
reporting methodologies; the data do not identify the specific reason for each disagreement. If the
SEC wants N-PORT to carry a comparable signal of issuer distress, a separate, explicitly defined field
for whether the issuer is in a pending bankruptcy case may serve better than redefining the
security-level default item, whose contractual meaning can legitimately differ across instruments. The
form should also say whether loans are covered. Until then, users should not read one fund's answer as
evidence about the issuer.

\paragraph{Use of AI.} Anthropic's Claude Opus 5.5 model was used to develop the analysis code;
extract and analyse the public SEC data; classify debtor status and read contract terms and other
public source documents; and generate the figures. The author wrote the first draft of the
manuscript, selected the research question, approved the research
protocol and its amendments, reviewed the evidence and conclusions, directed revisions, and accepts
responsibility for the final manuscript.

\clearpage
\appendix
\section*{Appendix A: contract sources}
\setcounter{table}{0}\renewcommand{\thetable}{A\arabic{table}}

\begin{table}[H]
\centering\footnotesize
\caption{Governing instruments behind Table~\ref{tab:contracts}. Accession numbers identify EDGAR filings; EoD: event-of-default
clause; Acc.: acceleration clause. Loan CUSIPs are matched by borrower, agreement and tranche, since
credit agreements do not print them. Links and the quoted clauses are in Appendix~B.}\label{tab:sources}
\begin{tabularx}{\textwidth}{@{}>{\raggedright\arraybackslash}p{2.9cm} >{\raggedright\arraybackslash}p{2.6cm} >{\raggedright\arraybackslash}X >{\raggedright\arraybackslash}p{2.1cm} >{\raggedright\arraybackslash}X@{}}
\toprule
Instrument (held CUSIP) & Obligor & Filing (exhibit) & EoD; Acc. & Obligor a debtor (source) \\
\midrule
Frontier 8.000\% 1L notes 2027 (35906ABC1) & Frontier Communications Corp & 0001140361-19-005134 (EX-4.1) & \S7.01(vi); \S7.02(c), automatic & Yes, 14 Apr 2020 (petition 8-K 0001140361-20-008876) \\
Frontier term B-1 loans (35906EAJ9) & Frontier Communications Corp & 0001193125-17-062225 (EX-10); joinder 0000020520-17-000018 (EX-10.1) & \S7.01(i); \S7.01, automatic & Yes, same source \\
Intelsat 8.50\% notes 2024 (45824TBA2) & Intelsat Jackson Holdings S.A. & 0001193125-18-277618 (EX-99.2) & \S6.01(e); \S6.02(b), automatic & Yes, 13 May 2020 (DIP agreement, 8-K 0001193125-20-171951) \\
Intelsat term loans, 2011 agreement (L5137LAF2) & Intelsat Jackson Holdings S.A. & 0001193125-14-060573 (EX-4.7, 2013 conformed text) & \S12.5; \S12, automatic & Yes, same source \\
Chesapeake 11.5\% 2L notes 2025 (165167DD6) & Chesapeake Energy Corp & 0001104659-19-076135 (EX-4.1) & \S6.01(6); \S6.02, automatic & Yes, 28 Jun 2020 (petition 8-K 0001104659-20-077745) \\
Chesapeake FLLO term loan (16516FAH2) & Chesapeake Energy Corp & 0001104659-19-076135 (EX-4.3) & \S10.5; \S10, automatic & Yes, same source \\
Endo 6.125\% 1L notes 2029 (29280BAA3) & Endo Luxembourg Finance Co I and Endo U.S. Inc.; parent Endo International plc & 0001193125-21-094652 (EX-4.1) & \S6.01(a)(7); \S6.02, automatic on a parent filing & Parent yes, 16 Aug 2022 (petition 8-K 0001193125-22-222687); Endo U.S. yes (plan, 0001193125-24-000733); Lux.\ Finance I not named on EDGAR \\
Endo 2021 term loans (L2968EAE2) & Endo Luxembourg Finance Co I; parent Endo International plc & 0001193125-21-094652 (EX-10.1) & \S7.01(i); \S7.01, automatic for parent or borrower & Parent yes, as above \\
Par Pharmaceutical 7.500\% 1L notes 2027 (69888XAA7) & Par Pharmaceutical, Inc.; parent Endo International plc & 0001193125-19-090665 (EX-4.1) & \S6.01(7); \S6.02, automatic on a parent filing & Yes (plan \S1.1.589, 0001193125-24-075864) \\
Whiting 6.625\% notes 2026 (966387BG6) & Whiting Petroleum Corp & 0001193125-17-380119 (EX-4.2) & \S6.01(9); \S6.02, automatic & Yes, 1 Apr 2020 (petition 8-K 0001193125-20-093737) \\
\bottomrule
\end{tabularx}
\end{table}

\section*{Appendix B: quoted clauses}
All links are to the SEC's EDGAR archive and were opened on 29 October 2026. ``EoD'' quotes the
bankruptcy event of default and ``Acc.'' the acceleration clause.

{\footnotesize
\begin{description}[leftmargin=0pt,itemsep=3pt]
\item[Frontier 8.000\% first-lien notes due 2027.] \url{https://www.sec.gov/Archives/edgar/data/20520/000114036119005134/s002637x3_ex4-1.htm}.
EoD \S7.01(vi): ``the Company shall commence a voluntary case under any applicable bankruptcy, insolvency or
other similar law''. Acc.\ \S7.02(c): ``If an Event of Default described in Section 7.01(vi) occurs and is
continuing, then the principal amount of all the Outstanding Notes, and the interest accrued thereon, if any,
shall become and be immediately due and payable without any declaration or other act on the part of the
Trustee or any Holder.''
\item[Frontier term B-1 loans.] Credit agreement \url{https://www.sec.gov/Archives/edgar/data/20520/000119312517062225/d356289dex10.htm};
increase joinder adding the B-1 tranche \url{https://www.sec.gov/Archives/edgar/data/20520/000002052017000018/ftr-20170630xex10_1.htm}.
EoD \S7.01(i): ``the Borrower or any of its Principal Subsidiaries shall (i) voluntarily commence any proceeding
or file any petition seeking liquidation, reorganization or other relief under any Federal or state
bankruptcy, insolvency, receivership or similar law''. Acc.: ``in case of any event with respect to the
Borrower described in clause (h) or (i) of this Section 7.01, the Commitments shall automatically terminate and
the principal of the Loans then outstanding \ldots{} shall automatically become due and payable''.
\item[Intelsat Jackson 8.50\% notes due 2024.] \url{https://www.sec.gov/Archives/edgar/data/1525773/000119312518277618/d627436dex992.htm}.
EoD \S6.01(e): ``the Issuer or any Significant Subsidiary pursuant to or within the meaning of any Bankruptcy
Law: (i) commences a voluntary case''. Acc.\ \S6.02(b): ``in the case of an Event of Default arising under
Section 6.01(e) or (f), all outstanding Notes will become due and payable without further action or notice.''
\item[Intelsat Jackson term loans (2011 agreement, 2013 conformed text).] \url{https://www.sec.gov/Archives/edgar/data/1525773/000119312514060573/d597626dex47.htm}.
EoD \S12.5: ``The Borrower or any Specified Subsidiary shall commence a voluntary case, proceeding or action
concerning itself under (a) Title 11 of the United States Code''. Acc., \S12: ``if an Event of Default
specified in Section 12.5 shall occur with respect to the Borrower or any Specified Subsidiary, the result
that would occur upon the giving of written notice by the Administrative Agent \ldots{} shall occur
automatically''. Debtor status: DIP agreement naming ``Intelsat Jackson Holdings S.A., as a Debtor and Debtor
in Possession'', \url{https://www.sec.gov/Archives/edgar/data/1525773/000119312520171951/d927580dex101.htm}.
\item[Chesapeake 11.5\% second-lien notes due 2025.] \url{https://www.sec.gov/Archives/edgar/data/895126/000110465919076135/tm1926440d1_ex4-1.htm}.
EoD \S6.01(6): ``any Obligor pursuant to or within the meaning of any Bankruptcy Law: (A) commences a voluntary
case or proceeding''. Acc.\ \S6.02: ``If an Event of Default specified in clauses (6) or (7) of Section 6.01
above occurs, all unpaid principal \ldots{} will become due and payable, without any declaration or other act
on the part of the Trustee or any Holder.''
\item[Chesapeake FLLO term loan.] \url{https://www.sec.gov/Archives/edgar/data/895126/000110465919076135/tm1926440d1_ex4-3.htm}.
EoD \S10.5: ``The Borrower or any Specified Subsidiary shall commence a voluntary case, proceeding or action
concerning itself under (a) Title 11''. Acc.: ``if an Event of Default specified in Section 10.5 shall occur
with respect to the Borrower, clauses (x) and (y) below shall occur automatically'', clause (y) making the
Loans ``forthwith due and payable''.
\item[Endo 6.125\% notes due 2029.] \url{https://www.sec.gov/Archives/edgar/data/1593034/000119312521094652/d120772dex41.htm}.
EoD \S6.01(a)(7): ``the Parent, the Issuers or any Significant Subsidiary: (a) commences a voluntary insolvency
proceeding''. Acc.\ \S6.02: ``If an Event of Default specified in Sections 6.01(a)(7) and 6.01(a)(8) hereof
with respect to the Parent shall occur, such amounts with respect to all the Notes shall become automatically
due and payable immediately without any declaration or further action or notice''.
\item[Endo 2021 term loans.] \url{https://www.sec.gov/Archives/edgar/data/1593034/000119312521094652/d120772dex101.htm}.
EoD \S7.01(i): ``Parent or any Material Subsidiary shall (i) voluntarily commence any proceeding or file any
petition seeking \ldots{} relief under any Federal, state or foreign bankruptcy \ldots{} law''. Acc.: ``in case of
any event with respect to Parent or any Borrower described in clause (h) or (i) of this Section 7.01, \ldots{}
the principal of the Loans then outstanding \ldots{} shall automatically become due and payable''.
\item[Par Pharmaceutical 7.500\% notes due 2027.] \url{https://www.sec.gov/Archives/edgar/data/1593034/000119312519090665/d694160dex41.htm}.
EoD \S6.01(7): ``the Parent, the Issuer or any Significant Subsidiary: (a) commences a voluntary insolvency
proceeding''. Acc.\ \S6.02: ``If an Event of Default specified in Sections 6.01(7) and 6.01(8) hereof with
respect to the Parent shall occur, such amounts with respect to all the Notes shall become automatically due
and payable immediately''; otherwise acceleration requires a declaration by the trustee or 25\% of holders.
Debtor status: plan \S1.1.589, \url{https://www.sec.gov/Archives/edgar/data/1593034/000119312524075864/d813107dex21.htm}.
Endo's 8-K of 19 August 2022 lists the 2021 credit agreement and both note issues as accelerated by the
filing: \url{https://www.sec.gov/Archives/edgar/data/1593034/000119312522225297/d389791d8k.htm}.
\item[Whiting 6.625\% notes due 2026.] \url{https://www.sec.gov/Archives/edgar/data/1255474/000119312517380119/d505619dex42.htm}.
EoD \S6.01(9): ``the Company, any Significant Subsidiary of the Company or any group of Subsidiaries \ldots{}
(a) commences a voluntary case''. Acc.\ \S6.02: ``if an Event of Default specified in clause (9) or (10) of
Section 6.01 hereof occurs with respect to the Company \ldots{} all outstanding Notes shall become due and
payable without further action or notice''.
\item[Counterexample: Valley National 6.25\% subordinated notes due 2032.] \url{https://www.sec.gov/Archives/edgar/data/714310/000119312522247735/d405024dex42.htm}.
EoD \S5.1(b): ``the Company or a Principal Subsidiary Bank shall commence a voluntary case under any applicable
bankruptcy, insolvency or other similar law''. Acc.\ \S5.2: ``the Trustee or the Holders of not less than
25\% in aggregate principal amount of the Outstanding Notes may declare the principal of all the Notes to be
due and payable immediately''.
\item[Petition 8-Ks.] Frontier \url{https://www.sec.gov/Archives/edgar/data/20520/000114036120008876/form8k.htm};
Intelsat \url{https://www.sec.gov/Archives/edgar/data/1525773/000119312520142170/d932199d8k.htm};
Chesapeake \url{https://www.sec.gov/Archives/edgar/data/895126/000110465920077745/tm2023599-1_8k.htm};
Endo \url{https://www.sec.gov/Archives/edgar/data/1593034/000119312522222687/d394120d8k.htm};
Whiting \url{https://www.sec.gov/Archives/edgar/data/1255474/000119312520093737/d899645d8k.htm}.
\end{description}
}

\bibliographystyle{plainnat}
\begin{thebibliography}{99}
\bibitem[Cavani et~al.(2025)Cavani, Maggiori, and Schreger]{cavani2025}
Cavani, F., M.~Maggiori, and J.~Schreger (2025). GCAP: Public security-level data on U.S.\ fund
holdings. NBER Working Paper 34494. \url{https://www.nber.org/papers/w34494}.
\bibitem[Cici et~al.(2011)Cici, Gibson, and Merrick]{cici2011}
Cici, G., S.~Gibson, and J.~J. Merrick~Jr. (2011). Missing the marks? Dispersion in corporate bond
valuations across mutual funds. \emph{Journal of Financial Economics} 101(1), 206-226.
\doi{10.1016/j.jfineco.2011.02.001}.
\bibitem[{Empirasign}(n.d.)]{empirasign}
Empirasign (n.d.). Surfacing outliers: N-PORT valuation outliers report.
\url{https://www.empirasign.com/blog/Surfacing-Outliers/}.
\bibitem[{ICE}(2026)]{ice_nport}
ICE (2026). ICE Regulatory Solutions: Form N-PORT service.
\url{https://www.ice.com/publicdocs/ICE_Regulatory_Solutions_Form_N_Port_Service.pdf}.
\bibitem[Jacobson and Kowara(2021)]{morningstar2021}
Jacobson, E., and M.~Kowara (2021). Bond pricing: Agreeing to disagree. Morningstar research paper.
\url{https://www.morningstar.com/content/dam/marketing/shared/research/foundational/Bond_Pricing_2021.pdf}.
\bibitem[Manoel(2026)]{manoel2026}
Manoel, P. (2026). Connections and creditor outcomes in Chapter 11. SSRN Working Paper 6226282.
\url{https://papers.ssrn.com/sol3/papers.cfm?abstract_id=6226282}.
\bibitem[George and Shin(2026)]{george2026}
George, L., and C.~Shin (2026). The rise of semiliquid funds for retail investors. FEDS Notes, Board of
Governors of the Federal Reserve System, 14 October 2026.
\url{https://www.federalreserve.gov/econres/notes/feds-notes/the-rise-of-semiliquid-funds-for-retail-investors-20261014.html}.
\bibitem[{SEC}(2016)]{sec_2016_release}
SEC (2016). Investment Company Reporting Modernization, Release Nos.\ 33-10231, IC-32314 (13 October
2016). \url{https://www.sec.gov/files/rules/final/2016/33-10231.pdf}.
\bibitem[{SEC}(2026)]{sec_2026_proposal}
SEC (2026). Form N-PORT reporting, proposed rule, Release No.\ IC-35962, 91 Fed.\ Reg.\ 8582 (23 February
2026), at 8595. \url{https://www.govinfo.gov/content/pkg/FR-2026-02-23/pdf/2026-03460.pdf}.
\bibitem[{SEC}(2020)]{corr_gs2020}
SEC (2020). Goldman Sachs Trust, correspondence, 28 August 2020.
\url{https://www.sec.gov/Archives/edgar/data/822977/000119312520234964/filename1.htm}.
\bibitem[{SEC}(2021)]{corr_rivernorth2022}
SEC (2021). RiverNorth Specialty Finance Corp., correspondence, 13 October 2021 (released on EDGAR 11 October 2022).
\url{https://www.sec.gov/Archives/edgar/data/1644771/000139834422020015/filename1.htm}.
\bibitem[{SEC}(n.d.a)]{sec_nport_form}
SEC (n.d.a). Form N-PORT, general instructions and Part C. \url{https://www.sec.gov/files/formn-port.pdf}.
\bibitem[{SEC}(n.d.b)]{sec_nport_data}
SEC (n.d.b). Form N-PORT data sets.
\url{https://www.sec.gov/data-research/sec-markets-data/form-n-port-data-sets}.
\bibitem[{US Code}(n.d.)]{usc362}
11 U.S.C.~\S362, Automatic stay. \url{https://www.law.cornell.edu/uscode/text/11/362}.
\bibitem[{US Courts}(n.d.)]{uscourts_ch11}
United States Courts (n.d.). Chapter 11: Bankruptcy basics.
\url{https://www.uscourts.gov/court-programs/bankruptcy/bankruptcy-basics/chapter-11-bankruptcy-basics}.
\bibitem[{Valley National Bancorp}(2022)]{valley2022}
Valley National Bancorp (2022). Second supplemental indenture, 6.25\% fixed-to-floating rate subordinated
notes due 2032, Form 8-K Exhibit 4.2, Sections 5.1 and 5.2.
\url{https://www.sec.gov/Archives/edgar/data/714310/000119312522247735/d405024dex42.htm}.
\end{thebibliography}

\end{document}
